命題3が抑えるのは種に混じるIDの\textit{割合}であり，未知が大半を占めるテスト集合では割合$q$の種がIDのほぼ全数になる（\secref{in200}）．そこで，step-upの臨界値を$c_i=\min(q\,i/m,\,\tau)$（Storey型では$q/\hat{\pi}_0$）に置き換えた規則も用意した．棄却された$p$はすべて$\tau$以下なので種に混じるIDの\textit{数}の期待値は$\tau m_0$以下であり，単純BH法では$p\le q\,i^*/m$（$i^*$は棄却数）も満たすので，PRDSのもとで自己整合な手続きとしてFDR$\le q\,m_0/m$が保たれる~\cite{blanchard2008two}．$\tau{=}0.05$は選ばずに固定した．この規則（$\tau{=}0.05$，$q{=}0.2$）は交差適合した$\sigma$について実行前に登録し，上限なしと同じ実行の中で評価した（付録\ref{sec:app_hyp}）．$\sR$との組合せは評価していない．

\textbf{ImageNet-1Kでは変わらない．}上限が効くのはStorey型の閾値が$\tau$を超えるSSB-hardとOpenImage-O（iNaturalistではわずか）で，SSB-hardでは種が折の\NCapSeedSsbHardStoreyRatio\%から\NCapSeedSsbHardStoreycRatio\%に減り，種に混じるIDの割合は\NCapSeedSsbHardStoreyIdfrac{}から\NCapSeedSsbHardStoreycIdfrac{}になる．near-OODの\FPR{}は単体で\NCapStoreycAloneNearFpr\%（上限なし\NAloneNearFpr\%，対応のある差\NCapStoreycAloneNearDfprVsStorey{}），$\times$TINSで\NCapStoreycNearFpr\%（+\NCapStoreycNearDfprVsStorey{}），far-OODとFour-OODの単体は\NCapStoreycAloneFarFpr\%，\NCapStoreycAloneFourFpr\%で，いずれも上限なしと0.25ポイント以内である．

\textbf{混入率が高いベンチマークでは害がほぼ消える．}ImageNet-200のnear-OODでは\FPR{}が\NInTwoZeroZeroMainNearFpr\%から\NInTwoZeroZeroStoreycNearFpr\%へ下がり（差\NInTwoZeroZeroStoreycNearDfprVsStorey{} [\NInTwoZeroZeroStoreycNearDfprVsStoreyLo, \NInTwoZeroZeroStoreycNearDfprVsStoreyHi]），交差適合した$\pP$だけ（\NInTwoZeroZeroLpNearFpr\%）より\NInTwoZeroZeroStoreycNearDfprVsLpAbs{}ポイント低い（[\NInTwoZeroZeroStoreycNearDfprVsLpLo, \NInTwoZeroZeroStoreycNearDfprVsLpHi]）．far-OODは\NInTwoZeroZeroStoreycFarFpr\%で交差適合した$\pP$だけ（\NInTwoZeroZeroLpFarFpr\%）と並ぶ．CIFAR-100ではnear-OODが\NCifarMainNearFpr\%から\NCifarStoreycNearFpr\%（交差適合した$\pP$だけ\NCifarLpNearFpr\%），far-OODが\NCifarMainFarFpr\%から\NCifarStoreycFarFpr\%になり，交差適合した$\pP$だけ（\NCifarLpFarFpr\%）との差は+\NCifarMainFarDfprVsLp{}から+\NCifarStoreycFarDfprVsLp{} [\NCifarStoreycFarDfprVsLpLo, \NCifarStoreycFarDfprVsLpHi]ポイントに縮む．\secref{rec}の格子の\NCaprecNCells{}条件でも，上限付きの規則が$\pP$だけより1ポイント以上悪い条件は混入率30\%以上の\NCaprecNHi{}条件で\NCaprecStoreycHiNWorseLpOne{}，$r\ge10$・混入率30\%未満の\NCaprecLoROneZeroN{}条件では上限なしとの差が最大\NCaprecLoROneZeroMaxStoreycMinusStorey{}ポイントと利得を保つ．

\textbf{位置づけ．}v18--v21では，上限付きの種がImageNet-1Kで上限なしと変わらず，混入率の高いベンチマークで$\pP$だけと同等以上になることから，混入率が未知の運用に上限付きを推していた．$\sR$と上限の組合せは評価していないので，本稿の推奨は結論のとおりとし，上限付きの規則は$\sigma$で登録した変種として報告する．以下の格子の値も$\sigma$のものである．
